\chapter{Simulation and Evaluation Methodology}

The evaluator is hierarchical: inexpensive checks reject invalid candidates
before costly transient and qualification work. \Cref{fig:optimization-flow}
shows the complete search and verification boundary.

\begin{figure}[p]
  \centering
  \includegraphics[width=\textwidth,height=0.82\textheight,keepaspectratio]{optimization-flow.pdf}
  \caption{Candidate search, selection, and verification flow. PPO changes
  bounded parameter-grid indices and receives metrics from the common staged
  evaluator. Exact cache hits reuse deterministic results. Candidate generation
  is followed by strict nominal filtering; optional PVT and final-fidelity
  checks occur only after a feasible design exists. Training saves a policy,
  while inference keeps policy weights fixed.}
  \label{fig:optimization-flow}
\end{figure}

\section{Evaluation stages}

\begin{table}[htbp]
\centering
\caption{Principal stages of a receiver evaluation.}
\label{tab:stages}
\begin{tabularx}{\textwidth}{@{}l X X@{}}
\toprule
Stage & Purpose & Representative outputs \\
\midrule
Validation & Check parameters, files, models, and conditions & Structured setup failure \\
DC & Establish bias, balance, common mode, and headroom & CTLE power and operating point \\
AC & Characterize differential transfer & Peaking, peak frequency, group delay \\
CTLE transient & Reject gross swing or rail violations & Input/output swing \\
Channel & Validate and characterize the S4P network & Loss, delay, passivity, causality \\
Noise & Integrate input-referred noise & RMS noise from \SI{10}{MHz}--\SI{5}{GHz} \\
HD3 & Fit the harmonic content of a sine response & Third-to-fundamental ratio \\
Receiver transient & Evaluate the channel-driven CTLE and DFE & Eye, margin, errors, transient power \\
\bottomrule
\end{tabularx}
\end{table}

\section{Fidelity levels}

Screening runs DC, AC, and a short channel-free diagnostic. Training adds the
channel and a 128-bit receiver evaluation. Candidate fidelity uses 512 bits and
enables noise and one-amplitude HD3. Final fidelity uses a 1024-bit request and
additional HD3 amplitudes. The number of useful post-warm-up bits can be smaller
than the requested sequence; result tables should therefore report the recorded
\code{simulated\_bits} value rather than infer it from the fidelity name.

\section{Caching and provenance}

The content-addressed cache is valid only when the complete evaluation identity
matches. The identity includes design parameters, environment conditions,
channel checksum and port map, model and bench checksums, schema versions, and
simulator configuration. Retryable execution failures are not treated as valid
cached circuit results.
